\documentclass[11pt,a4paper]{article}
\usepackage[T1]{fontenc}
\usepackage[utf8]{inputenc}
\usepackage{lmodern}
\usepackage[margin=1.65cm,top=1.3cm,bottom=1.1cm]{geometry}
\usepackage{booktabs}
\usepackage{graphicx}
\usepackage{xcolor}
\usepackage{microtype}
\usepackage{amsmath}
\usepackage{caption}
\captionsetup{font=footnotesize,labelfont=bf}
\pagestyle{empty}
\setlength{\parskip}{0.30em}
\setlength{\parindent}{0pt}

\newcommand{\hd}[1]{\vspace{0.25em}\textbf{\large #1}\par\vspace{0.1em}}

\begin{document}
\small

{\Large\bfseries Sea-ice sub-cycling in the FESOM2 GPU port: what pays}\\[0.2em]
{\small N. Koldunov, 6 August 2026 --- summary of the mEVP communication study}

\vspace{0.5em}\hrule\vspace{0.7em}

\hd{The one change that pays}

\textbf{Exchange the ice-velocity halo every 4th mEVP sub-cycle instead of after every one.}
mEVP runs 120 sub-cycles per model time step and currently exchanges after every one; this
exchanges 30 times instead of 120 and lets the halo be up to three sub-cycles old. It is one
conditional, switched on at run time, and it is an \emph{approximation} --- which is why most of
the work went into measuring what it costs.

\hd{What it buys}

Each mesh at a node count where the model still scales --- not the largest available, which would
flatter the result. Change of the whole model time step; GPU, same allocation, faster of two
repetitions.

\begin{center}\small
\begin{tabular}{@{}llrrrr@{}}
\toprule
\textbf{mesh} & \textbf{operating point} & \textbf{standard} & \textbf{every 4th} & \textbf{every 8th} \\
& & s/step & sub-cycle & sub-cycle \\
\midrule
CORE2 (127 k nodes) & 1 node (4 GPUs), knee at 2 & 0.0666 & $-12.6\%$ & $-14.4\%$ \\
fArc (638 k) & 4 nodes (16 GPUs), flat past 4 & 0.1173 & $-11.4\%$ & $-14.3\%$ \\
DARS (3.2 M) & 8 nodes (32 GPUs), still scaling & 0.1720 & $-14.2\%$ & $-15.5\%$ \\
NG5 (7.4 M) & 16 nodes (64 GPUs), still scaling & 0.2424 & $-8.5\%$ & $-9.9\%$ \\
\bottomrule
\end{tabular}
\end{center}

\textbf{So: 11--14\% off a production GPU run at $K=4$ on three of the four meshes (8.5\% on NG5),
and it is not an artefact of running past the scaling limit} --- each row is that mesh's own
sensible operating point. (Run past the limit the number is larger --- $-25.5\%$ for CORE2 on
16 nodes --- but nobody runs CORE2 on 16 nodes, and that number is not the one to quote.)

\textbf{Caveat: this is a GPU result.} On CPU the same change is worth $-1$ to $-3.4\%$, because a
CPU process holds far less of the mesh and the ice exchange is a much smaller share of the step.

\hd{What it costs}

One-year CORE2 integrations on the CPU backend, which is exactly reproducible, so the whole
difference is the approximation. Yardstick: two independent and both correct implementations of
the same mEVP scheme --- our C port and Fortran FESOM2 --- over the same year.

\begin{center}\small
\begin{tabular}{@{}lccccc@{}}
\toprule
pattern correlation, 1 year & SST / SSS / SSH & ice concentration & ice thickness & $u_{\rm ice}$ & $v_{\rm ice}$ \\
\midrule
C port versus Fortran (yardstick) & 1.00000 & 0.99999 & 0.99998 & 0.95438 & 0.93908 \\
\textbf{every 4th sub-cycle} & \textbf{1.000000} & \textbf{0.999995} & \textbf{0.999994} & \textbf{0.963328} & \textbf{0.950003} \\
every 8th sub-cycle & 1.000000 & 0.999992 & 0.999981 & 0.961359 & 0.944857 \\
\bottomrule
\end{tabular}
\end{center}

Every field stays inside the yardstick at $K\le8$. Ice velocity is the sensitive field and the
margin there is only 1.1--1.5$\times$, so smaller $K$ is the defensible choice --- the error is
monotone in $K$, and $K=4$ keeps about 87\% of the speed of $K=8$. Integrated quantities are
untouched: Northern Hemisphere ice extent deviates by 0.23\% of its seasonal amplitude at $K=8$,
and surface temperature, salinity and sea surface height correlate at 1.000000.

\textbf{Not yet measured: accuracy at high resolution.} The accuracy test is CORE2; the speed
results are on fArc, DARS and NG5. $K=4$ is recommended \emph{scoped to CORE2} until a
high-resolution year has been run.

\hd{What does not pay --- measured, so we do not have to argue about it}

\textbf{Danilov's divergence-form reformulation} (carry $\nabla\!\cdot\!\sigma$ at nodes instead of
$\sigma$ at elements). Exact --- it reproduces $\nabla\!\cdot\!\sigma$ to $10^{-15}$ relative ---
but worth $-0.3$ to $+0.1\%$ on CPU and $0$ to $+1.6\%$ on GPU. The sub-cycle is limited by
communication and by the number of GPU operations issued, not by arithmetic or memory traffic, so
removing element arrays removes work that was not costing anything: a control that deleted three
element arrays outright changed the time by exactly 0.00\%.

\textbf{The exact wide halo} --- Danilov's double-halo proposal (exchange every $K$-th sub-cycle
but recompute a $K$-deep ghost ring, so the answer is unchanged bit for bit). Certified exact, and
worth $-1.8$ to $-6.8\%$ at the operating points above, against the approximation's $-8.5$ to
$-15.5\%$. It needs $K=8$: at $K=2$ it \emph{costs} time, because the duplicated ring is not
amortised. On CPU it never wins on any mesh --- at 2048 processes on fArc the ring a process must
advance is larger than its own sub-domain. It remains the right choice for anyone who needs the
answer unchanged bit for bit.

\hd{Recommendation}

Adopt the delayed exchange at $K=4$ as an opt-in option for GPU production runs, after a
high-resolution accuracy year. Keep the exact wide halo for anyone who needs the answer unchanged
bit for bit. The divergence form is correct and elegant, but it is not a speed lever on this
hardware.

\end{document}
